% 第8回　店舗設計・空間構成
\session{8}{2}{11月18日（水）}{店舗設計・空間構成}{AIと空間記述}{yes}

\goals{
  \item 空間をゾーンと要素（素材、光、音、サイン）に分けて記述できる
  \item 言葉による空間記述を、画像生成のプロンプトに変換できる
  \item 生成画像を見て、コンセプトとの一致点とずれを言語化できる
  \item 生成画像の扱い（イメージであること、著作権と商標への注意）を理解する
}

\fillin{コンセプト文}

\work[90mm]{ワーク1}{空間の要件と、ゾーンごとの記述}
\begin{promptbox}[空間を言葉にする]
次のコンセプトとペルソナをもとに、店舗の空間を記述してください。入口から順に、来店客が通るゾーンごとに、広さの印象、床と壁の素材、光の質、音、配置されている家具や展示物、スタッフとの距離感を具体的に書いてください。コンセプト：（コンセプト文）　ペルソナ：（要約）
\end{promptbox}
\begin{wstabh}{colspec={Q[l,wd=16mm,m] X[l,m] X[l,m] X[l,m] X[l,m]},row{2-Z}={ht=15mm}}
  ゾーン & 広さ・素材 & 光・音 & 家具・展示物 & スタッフとの距離 \\
  受付 & & & & \\ 展示 & & & & \\ 商談 & & & & \\ 待合 & & & & \\ 整備 & & & & \\ & & & & \\
\end{wstabh}

\work[60mm]{ワーク2}{画像生成のプロンプトに変換し、条件を変えて生成する}
\begin{promptbox}[画像生成のプロンプトに変換する]
上の空間記述のうち「（ゾーン名）」の部分を、画像生成AI向けのプロンプトに変換してください。視点（入口から見た、俯瞰など）、時間帯、素材、光、人の有無を含め、英語と日本語の両方で出してください。
\end{promptbox}
\begin{wstabh}{colspec={Q[c,wd=6mm,m] Q[l,wd=16mm,m] X[2,l,m] X[2,l,m] X[2,l,m]},row{2-Z}={ht=14mm}}
  & ゾーン & 変えた条件（視点・時間帯など） & 表現したかったこと & うまくいかなかったこと \\
  1 & & & & \\ 2 & & & & \\ 3 & & & & \\
\end{wstabh}

\work{ワーク3}{画像を並べて言葉にする}
\twoboxes{コンセプトと一致している点}{違和感・ずれ}{30mm}
\answerbox[空間記述の修正点]{20mm}

\work[90mm]{スケッチ}{平面のゾーニング（手描き可。入口・動線・各ゾーンを書き込む）}
\gridbox{85mm}

\begin{caution}{生成画像を使うときの注意}
□ 実在の店舗、ブランドのロゴ、既存の作品に似た画像は使わない
□ 提案資料に載せる場合は「AIによる生成イメージ」と明記する
□ 画像はあくまでイメージ。寸法や法規の裏づけは別に必要
\end{caution}

\begin{nexttime}
  \item 空間コンセプト図（ゾーニングを示した平面スケッチ。手描き可）を仕上げる
  \item 生成画像2〜3点と、それぞれで表現したかったこと・うまくいかなかったことのメモ（上のワーク2）
\end{nexttime}
\aimemo
